\section{AI Demand、制造反向规划与 Warp Speed}

讲者认为 AI 讨论过度集中在 supply，也就是模型、compute 和 benchmark。电力本身创造条件，真正形成生产率的是围绕电力设计的新机器和流程；类似地，模型能力需要 data integration、ontology、工具、workflow 和 outcome feedback 才能改变组织决策。这个 demand-side 观点与本讲前半的 Apollo/Rubix 一致：基础模型并不会自动获得部署、权限和闭环。

\subsection{AI supply 之后的 Machine Layer}

Machine layer 将模型接到真实 entity、工具和责任体系。它决定模型看到哪些状态、能执行什么 action、何时需要审批、怎样处理失败。医疗、制造和公共服务的问题在 ChatGPT 之前已经存在；AI 改变的是可用能力和成本，不会自动定义问题、数据权利和成功标准。

\begin{figure}[H]
\centering
\includegraphics[width=0.94\textwidth]{images/10-ai-supply-demand.png}
\caption{AI 价值链：模型供给需要 machine/workflow layer 才能转成可验证 outcome。}
\figsource{依据字幕 30:00--33:30 重绘，`images/10-ai-supply-demand.png`。}
\end{figure}

\begin{knowledgebox}{读图：Machine layer 不是简单 RAG 包装}
它至少包括可靠数据管道、entity resolution、权限、工具执行、业务约束、人工协作和结果评估。一个回答正确率很高的模型，如果不能访问实时库存、不能尊重订单优先级、也不能把执行结果写回系统，就没有形成工业决策能力。
\end{knowledgebox}

\subsection{供应冲击为何要求反向规划}

传统 manufacturing planning 常从 customer orders 展开 bill of materials，再生成采购和生产计划。讲者用 BMW 的 wiring harness 供应波动举例：当关键小部件无法按时足量交付，系统必须从现有 inventory 反向计算哪些配置可生产、哪些订单价值更高、怎样调整工位和交付。旧软件若只支持正常方向，员工就会逃到 Excel 手工求解。

\begin{figure}[H]
\centering
\includegraphics[width=0.94\textwidth]{images/11-manufacturing-reverse-plan.png}
\caption{供应中断时，规划方向从订单→物料转为库存→可行配置→订单选择。}
\figsource{依据字幕 33:30--38:00 重绘，`images/11-manufacturing-reverse-plan.png`。}
\end{figure}

\begin{importantbox}{术语消化：Bill of Materials 与 feasible plan}
Bill of Materials，简称 BOM，是制造一个产品所需部件、数量和层级关系。Feasible plan 是在库存、设备、人员、交期和政策约束下真正可执行的方案。利润最高的订单若缺关键部件就不可行；部件齐全的方案若违反认证或产线能力也不可行。
\end{importantbox}

\subsection{Warp Speed 与制造软件批判}

Sankar 将 Warp Speed 描述为把长期工业项目中的 supply、production 和 workflow abstraction 产品化的平台。他批评 legacy manufacturing stack 在 COVID 和战争冲击中暴露出静态规划与系统割裂，并强调工程师要在 factory floor 观察真实工作。这里应把公司产品主张与通用机制分开：动态 constraint solving、设计—供应—生产反馈和可执行 workflow 是可迁移需求，具体产品效果需要客户侧指标验证。

\begin{warningbox}{讲者立场：产业竞争不是技术事实}
访谈把制造能力、国防供应链和国家竞争放在强烈战略叙事中。这些主张需要经济、历史和政策材料另行验证。讲义保留其对软件栈和 feedback 的系统观察，不把具体国家比较、产能数字或政策建议包装成由架构图自动推出的结论。
\end{warningbox}

\begin{knowledgebox}{课堂提示：为什么 factory-floor Excel 是强信号}
当一线人员把关键决策搬到 spreadsheet，通常说明正式系统缺少数据、速度、表达能力或信任。禁止 Excel 只会隐藏问题。更好的做法是观察 spreadsheet 中的 entity、constraint 和公式，识别哪些应该进入共享 ontology 与可审计 workflow。
\end{knowledgebox}

\subsection{本章小结}

AI 的价值来自 demand-side machine layer，制造案例则展示真实约束如何迫使规划方向改变。Warp Speed 的通用启示是把设计、供应、生产和反馈连接起来；产品主张和国家战略判断仍需独立证据。最后回到高风险软件的伦理问题：如何在隐私和安全之间获得更多能力，而不是只做零和取舍。

